\section{Introduction}
Battery-powered three-wheelers (easy-bikes, Mishuks, e-rickshaws) have become a mass transport mode in Bangladesh. A CLEAN/BWGED study estimates about 1.4~million easy-bikes consuming roughly 4{,}905~GWh of electricity per year \citep{clean2024}, whereas other sources put the number of electric three-wheelers at 3--4~million \citep{pureearth_smep}; we carry both figures as scenarios. Lead-acid batteries dominate \citep{mordor2026} and are short-lived: a battery set costing more than BDT~72{,}000 is reported to last 8--11~months \citep{fe2024}. A 2026 field study that surveyed 150 drivers in Gopalganj and Khulna documents prolonged hub charging, frequent overloading and lead-acid battery degradation \citep{chowdhury2026}.

This short life has three consequences. (i)~\emph{Economic}: replacement dominates operator cost. (ii)~\emph{Energy and carbon}: internal resistance and charger losses waste electricity, and each replacement carries manufacturing emissions. (iii)~\emph{Public health}: used lead-acid batteries (ULABs) are broken and smelted largely by the informal sector, a primary contributor to lead-pollution hotspots \citep{esdo2021}. The national public-health bulletin of IEDCR summarises the burden: about 36~million Bangladeshi children (60\%) are poisoned by lead, lead causes an estimated loss of about 20~million IQ points and 140{,}000 cardiovascular deaths, and the 2019 economic burden was US\$28.6~billion, 6--9\% of GDP \citep{das2026}. Pure Earth reports about 125~kg of lead in one easy-bike's battery set and about 167{,}000~t of lead waste per year from the electric three-wheeler sector \citep{pureearth_smep}.

LFP and NMC batteries offer higher round-trip efficiency, much longer cycle life and no lead, but cost more and are imported. Bangladesh is also hot: monthly mean 2\,m temperature in the study area ranges from \Tmin\ to \Tmax$^\circ$C (\cref{fig:temp}), and a 42$^\circ$C day has been reported in Chuadanga \citep{heat2025}. Heat accelerates ageing of every chemistry, but not equally, so whether lithium is preferable under local conditions, and by how much, is an empirical question.

\paragraph{Research gap.} Existing work compares lead-acid and lithium-ion for stationary storage \citep{keshan2016,yudhistira2022,jasper2025}, studies cell ageing \citep{yagci2026,jalkanen2015}, thermal safety \citep{schoberl2024}, charging infrastructure \citep{hasan2020}, or easy-bike operation and energy use for lead-acid only \citep{chowdhury2026,peakdemand}. We are not aware of a study that links satellite-reanalysis temperature to chemistry-specific ageing, replacement-aware cost, carbon and lead-flow accounting for Bangladesh easy-bikes in one transparent model. \todo{state the databases, search terms and dates of your final literature search}

\paragraph{Contributions.}
\begin{enumerate}[nosep,leftmargin=*]
\item A reproducible MATLAB/Simulink model (analytic and block-diagram versions cross-checked) of three chemistries driven by NASA POWER temperature.
\item A parameter matrix in which every value carries a published source or an explicit reason for being a scenario (\cref{tab:params}).
\item A carbon and lead-flow account for the fleet, with fleet-size and emission-factor sensitivity.
\item A four-scenario policy comparison, a break-even analysis, and a worst-case test of when LFP stops being cheaper.
\end{enumerate}
